% !TEX root = ../../proposal.tex

\section{Disentangled Representation Learning}
\label{sec:rw_disentanglement}

The goal of learning a representation in which each latent dimension corresponds to one
interpretable factor of variation was formalised by the disentanglement literature.
$\beta$-VAE~\citep{Higgins2016betaVAELB} showed that strengthening the VAE's KL bottleneck
encourages factorial posteriors at the cost of reconstruction quality. FactorVAE~\citep{disentangling_by_factorising}
refined this by penalising total correlation directly, producing cleaner factor separation
at comparable reconstruction. The DCI framework~\citep{eastwood2018framework} provided a
principled evaluation methodology: factor importance matrices derived from supervised
regressors yield per-dimension Disentanglement, per-factor Completeness, and overall
Informativeness scores that are independent of model architecture. Locatello et
al.~\citep{locatello2019challenging} delivered a decisive negative result: unsupervised
disentanglement is unidentifiable without inductive bias, and random seeds account for as
much variance in disentanglement scores as method choice. This finding directly motivates
the supervised regimes in Phase~2 --- Regime~C2 adds label-based factor supervision, and
Regime~D adds structural partitioning --- as the inductive biases needed to make
disentanglement reliable under correlated training data.

Traeuble et al.~\citep{traeuble2021disentangled} showed that when training factors are
correlated, standard independence-based objectives fail to disentangle them because the
model exploits the correlation as a prediction shortcut. Weak supervision over the factor
labels partially restores separation. This is the precise failure mode that Phase~2
Regime~B exhibits: correlation in the dSprites training distribution entangles orientation
and scale, causing held-out composition to degrade. The Regime~A--B--C1--C2--D ladder
in Phase~2 is directly motivated by this result. Roth et al.~\citep{roth2023disentanglement}
extended this analysis to show that data augmentation strategies that break the spurious
correlation can also help, complementing the supervision approach.

The adversarial mechanism used in Phase~2 for wrong-head suppression traces to Ganin and
Lempitsky~\citep{ganin2015unsupervised, ganin2016domain}, who introduced the Gradient
Reversal Layer (GRL) for domain-invariant feature learning in DANN. The GRL is identity
in the forward pass but negates the gradient in the backward pass, turning any attached
classifier into an adversary that instructs the encoder to remove a nuisance attribute
from its representation. Scaramuzzino et al.~\citep{attribute_disentanglement} applied
this mechanism directly to attribute disentanglement in fashion retrieval: a correct-head
classifier encourages each attribute to be encoded in its designated latent region, while
a GRL-coupled wrong-head classifier actively suppresses that attribute from all other
regions. Phase~2's SP-VAE design follows this pattern precisely, extending it to a
three-way latent partition with two independently supervised and adversarially excluded
factors. SepVAE~\citep{louiset2024sepvae} provides the structural partitioning inspiration:
by allocating common and salient variation to separate latent subspaces rather than
regularising within a shared code, it achieves more reliable pathology separation in
contrastive medical imaging datasets than independence penalties alone. SP-VAE adapts
the SepVAE structural insight to the single-dataset, multi-factor recombination setting
of correlated dSprites and combines it with GRL-based exclusion.
